% ============================================================
% §2  Related Work
% ============================================================

\section{Related Work}
\label{sec:related}

\subsection{Multivariate Time Series Anomaly Detection}
\label{sec:related_mtsad}

Deep learning approaches to unsupervised MTSAD fall into several well-defined families: reconstruction-based methods flag inputs that an encoder--decoder trained to reproduce normal data reconstructs poorly~\cite{zong2018dagmm,su2019omnianomaly,audibert2020usad,song2023memto}; prediction-based methods score deviations from a forecast of the expected next state~\cite{deng2021gdn}.
A more recent strand scores the discrepancy between learned and observed association structure, via transformer-learned temporal dependencies~\cite{xu2022anomalytransformer} or multi-scale contrastive views~\cite{yang2023dcdetector}; frequency-domain reconstruction has been extended with explicit channel-correlation discovery~\cite{wu2025catch}.
General-purpose time-series backbones and auxiliary training objectives have also been applied directly to MTSAD~\cite{tuli2022tranad,wu2023timesnet}.

All of these families, however, treat the training data as predominantly or entirely normal: when the training stream contains confirmed anomalous events (the contaminated setting of Section~\ref{sec:problem_formulation}), labeled information is either discarded or treated as noise that corrupts the learned normality model~\cite{wang2025nrdetector}.

\subsection{Label-Informed Anomaly Detection: Semi-supervised, PU, and Weakly Supervised}
\label{sec:related_semi}

Positive and Unlabeled (PU) learning formalizes learning from confirmed positive examples plus unlabeled data that may contain additional positives~\cite{bekker2020pusurvey,duplessis2014pu}; established solutions include non-negative risk estimators~\cite{kiryo2017nnpu}, class-prior-based probability correction~\cite{elkan2008pu}, and two-step techniques that extract reliable negatives before training a classifier~\cite{bekker2020pusurvey}.
Outside time series, these techniques have been adapted to anomaly detection through deviation networks with scarce labeled anomalies~\cite{pang2019devnet} and deep semi-supervised objectives~\cite{ruff2020deepsad}.

In the time-series domain, methods that incorporate anomaly labels into the representation learning objective itself are rare~\cite{wang2025nrdetector}.
A weakly supervised strand trains models to classify or rank windows from coarse segment-level annotations~\cite{sultani2018deepmil,lee2021wetas,liu2024treemil}, treating the label as the sole supervision signal, without any self-supervised reconstruction pretext.
Our use of labels differs in kind: rather than acting as the target of a classification or ranking objective, labels here enter the gradient of the masked-reconstruction pretext task, shaping what the encoder itself learns to represent.
Two earlier semi-supervised models address label-scarce multivariate time series: an autoregressive normality model with discriminative loss components that separate normal data from the few labeled anomalies~\cite{xue2022fewpositive}, and a variational autoencoder coupled with an active-learning labeling loop~\cite{huang2022slavae}.
In both, labels act through loss terms attached to a generative or predictive normality objective; neither employs a masked-reconstruction self-distillation pretext nor adversarial gradient-level suppression of anomaly information.
In the transfer setting, DACAD~\cite{darban2024dacad} exploits labeled anomalies from a related source domain through supervised contrastive learning; in our setting, the scarce labels reside in the target training stream itself.

The closest precedent is NRdetector~\cite{wang2025nrdetector}, which formulates point-level detection under noisy segment-level labels as a PU problem and identifies this as a novel setting that prior time-series anomaly detection methods only partially address.
Its framework is a pipeline: a pre-trained WETAS-derived backbone extracts temporal embeddings, and a separate PU classifier is trained on those fixed representations; the label signal guides the classifier's output, not the encoder's gradient.
Our approach differs along this axis: labeled anomaly information enters the gradient of the encoder during training itself, through three orthogonal mechanisms that shape what the model learns to represent rather than what it predicts.
To our knowledge, CSMAD is the first end-to-end MTSAD model that integrates labeled anomalies into the gradient of a masked-reconstruction self-distillation objective through gradient reversal.

\subsection{Masked Autoencoders and Self-Distillation in Anomaly Detection}
\label{sec:related_mae}

The masked autoencoder (MAE) of He et al.~\cite{he2022mae} demonstrates that masking random patches and reconstructing the missing regions produces strong transferable representations.
Our patch-based masking adapts this spatial-domain paradigm to windows of multivariate sensor channels; similar masking-based reconstruction objectives in some time-series models~\cite{fang2024tfmae} constitute independent developments, whereas our design follows directly from vision MAE.

Knowledge distillation has been applied to anomaly detection through teacher--student frameworks: a student trained to match a pre-trained teacher's representations fails to do so on anomalous inputs, revealing the anomaly as a representation gap~\cite{bergmann2020uninformed,deng2022reverse}.
A more self-contained formulation is self-distillation~\cite{zhang2022selfdistill}, which performs the distillation within a single architecture.
Ristea et al.~\cite{ristea2024sdmae} adapt this design to video anomaly detection, pairing a capacity-limited student decoder with a deeper teacher inside a masked autoencoder and using their output discrepancy as the anomaly score at inference.

We adapt this asymmetric teacher--student masked autoencoder design to multivariate time series, embedding it within the contaminated semi-supervised framework of Section~\ref{sec:problem_formulation}, where labeled anomalies guide training through the three pathways of Sections~\ref{sec:masking}--\ref{sec:label_training}.\footnote{%
  The self-distillation terminology is adopted from Zhang et al.~\cite{zhang2022selfdistill} and Ristea et al.~\cite{ristea2024sdmae}.
  We capitalize Teacher and Student when referring to CSMAD's own decoder branches and use lowercase for the general teacher--student paradigm.
  Unlike SDMAE, whose student decoder branches off from within the teacher decoder after its first transformer block, our Teacher and Student decoders are independent parallel branches off the shared encoder.
  The gradient reversal layer that adversarially suppresses anomaly information in the Student is absent from Ristea et al.~\cite{ristea2024sdmae}, who train without real labeled anomalies.
  The distinction between operating in the target/loss space and in the gradient space of the representation is elaborated in Section~\ref{sec:label_training}.}
